\chapter{\abstractname}

Engineering systems often describe one product with several models that share
information, so a change to one model has to be propagated to all other models
that depend on it. When many updates arrive concurrently, this propagation can
either be done inside each transaction (ACID), which keeps the models consistent
but limits throughput, or be moved to the background (BASE), which achieves high
throughput but leaves the models temporarily inconsistent without giving the
administrator control over how much inconsistency is accepted.

This thesis addresses this trade-off with a hybrid execution mode that decides
for each transaction whether its changes are propagated immediately or in the
background, so that both kinds of work can run side by side on the same data
without conflicting. The decision is made by a mode oracle, and this work
presents two of them, one that applies immediate propagation to a fixed share of
transactions and one that adapts to the current level of inconsistency in the
system. The hybrid mode is implemented in an existing simulator of multi-model
systems, which makes it possible to evaluate how well it balances consistency
and throughput.

The evaluation shows that the hybrid mode reaches a higher
throughput than pure ACID and, unlike pure BASE, keeps the inconsistency close
to a configured threshold. Under a static load, however, the adaptive oracle
gives no performance gain over the fixed share.


\makeatletter
\ifthenelse{\pdf@strcmp{\languagename}{english}=0}
{\renewcommand{\abstractname}{Kurzfassung}}
{\renewcommand{\abstractname}{Abstract}}
\makeatother

\chapter{\abstractname}

\begin{otherlanguage}{ngerman}
  In der Systementwicklung wird ein Produkt oft durch mehrere Modelle
  beschrieben, die Informationen teilen, sodass eine Änderung an einem Modell
  auf alle anderen Modelle übertragen werden muss, die davon abhängen. Wenn
  viele Änderungen nebenläufig eintreffen, kann diese Übertragung entweder
  innerhalb jeder Transaktion erfolgen (ACID), was die Modelle konsistent hält,
  aber den Durchsatz begrenzt, oder in den Hintergrund verlagert werden (BASE),
  was einen hohen Durchsatz erreicht, die Modelle jedoch zeitweise inkonsistent
  lässt, ohne dass der Administrator steuern kann, wie viel Inkonsistenz
  akzeptiert wird.

  Diese Arbeit begegnet diesem Zielkonflikt mit einem hybriden
  Ausführungsmodus, der für jede Transaktion entscheidet, ob ihre Änderungen
  sofort oder im Hintergrund übertragen werden, sodass beide Arten von Arbeit
  ohne Konflikte nebeneinander auf denselben Daten laufen können. Die
  Entscheidung trifft ein Modus-Orakel, von denen diese Arbeit zwei vorstellt,
  eines, das einen festen Anteil der Transaktionen sofort überträgt, und eines,
  das sich an den aktuellen Grad der Inkonsistenz im System anpasst. Der hybride
  Modus ist in einem bestehenden Simulator für Multi-Modell-Systeme
  implementiert, wodurch sich bewerten lässt, wie gut er Konsistenz und
  Durchsatz ausbalanciert.

  Die Evaluation zeigt, dass der hybride Modus einen
  höheren Durchsatz als reines ACID erreicht und, anders als reines BASE, die
  Inkonsistenz nahe an einem konfigurierten Schwellenwert hält. Unter statischer
  Last bringt das adaptive Orakel jedoch keinen Leistungsgewinn gegenüber dem
  festen Anteil.
\end{otherlanguage}

% Undo the name switch
\makeatletter
\ifthenelse{\pdf@strcmp{\languagename}{english}=0}
{\renewcommand{\abstractname}{Abstract}}
{\renewcommand{\abstractname}{Kurzfassung}}
\makeatother
